%From Jose Pinheiro <jcp@rice.research.bell-labs.com>
\documentclass[a4,semhelv,semcolor,portrait,11pt]{seminar}
%May want to add slidesonly after 11pt if using notes
%Bibliography may work with slidesonly if put it inside a slide
\centerslidesfalse
%\usepackage{S}
\usepackage{graphicx}
%\usepackage{pstricks}


%\newpagestyle{intro}%
%{\sl OVERALL TITLE \hfil SHORT TITLE}{\hfil \thepage}
%\pagestyle{intro}

\slideframe{none}

\graphicspath{{figs/}}
% \DeclareGraphicsExtensions{.eps.gz,.eps,.epsi.gz,.epsi,.ps,.ps.gz}
% \DeclareGraphicsRule{*}{eps}{*}{}
% \DeclareGraphicsRule{.epsi.gz}{eps}{.bb}{`zcat #1}
% \DeclareGraphicsRule{.eps.gz}{eps}{.bb}{`zcat #1}
% \DeclareGraphicsRule{.ps.gz}{eps}{.bb}{`zcat #1}
\setkeys{Gin}{width=\textwidth}

%Example: \fig{myfig}{An example of a good fit on worthless data}
%Inserts postscript figure myfig.ps into document, with caption given
%This is for .ps files not generated by transfig
%\fig{label=.ps prefix}{caption}
\newcommand{\fig}[2]{\begin{figure}[htbp]           %
    \begin{center}  
    % can also try \centerline{\includegraphics{}} or \centering b4 includeg.
   \leavevmode
   \includegraphics{#1}
%   \includegraphics[width=0.9\textwidth]{#1}
   \end{center}
    \caption{#2}
    \label{#1}
   \end{figure}}

\newcommand{\bs}[1]{\begin{slide*}\Ti{#1}}
\newcommand{\es}{\end{slide*}}

\newcommand{\bi}{\begin{itemize}}
\newcommand{\ei}{\end{itemize}}
\newcommand{\be}{\begin{enumerate}}
\newcommand{\ee}{\end{enumerate}}

\newcommand{\bm}[1]{{\boldsymbol {#1}}}
\newcommand{\Ti}[1]{\hfill {\textbf{\large \blue  #1}\\}%
  \vspace*{-1cm} \rule{\linewidth}{0.2mm}}

\newcommand{\mg}[1]{{\magenta #1}}
\newcommand{\rd}[1]{{\red #1}}
\newcommand{\bu}[1]{{\blue #1}}
\newcommand{\gr}[1]{{\green #1}}

\newcommand{\Cf}[1]{\textsf{#1}}
\newcommand{\splus}{{S-P\sc{lus}} \mbox{}}
\newcommand{\spc}{$\mbox{}$}
\newcommand{\sct}[1]{{\scriptsize #1}}
\newcommand{\sctm}[1]{{\scriptsize $#1$}}
\begin{document}

\vspace*{\stretch{1}}
\noindent\rule{\linewidth}{1mm}
\begin{flushright}
\huge
TITLE\\[2ex]
\Large Frank E Harrell Jr \\
\normalsize
Division of Biostatistics and Epidemiology \\
Department of Health Evaluation Sciences \\
University of Virginia School of Medicine \\
Box 800717 Charlottesville VA 22908 USA \\
{\tt fharrell@virginia.edu} \\
{\tt hesweb1.med.virginia.edu/biostat} \\
\end{flushright}
\rule{\linewidth}{1mm}
\vspace*{\stretch{2}}
\begin{center}
\textsc{MEETING \\ DATE}
\end{center}
\thispagestyle{empty}

%Shortcut for below:
\bs{Outline}
\be
\item ....
   \bi
   \item ....
   \ei
\ee
\es

\begin{slide*}
  \Ti{Outline}
  \begin{enumerate}
    \item \rd{Introduction}
    \begin{itemize}
    \item motivating examples of \mg{linear} and \mg{non-linear}
      mixed-effects models 
   \item   $x^{2} - \sum_{i=1}^{n}\chi^{2}$
    \item a sample of \splus commands for \mg{fitting} and
      \mg{analyzing} mixed-effects models. 
    \end{itemize}
  \item \rd{Grouped data}
    \begin{itemize}
    \item \spc \mg{repeated measures}, \mg{longitudinal} data, \mg{growth
        curve} data 
    \item  \spc \mg{multiply-nested} data
    \item Data structures, \mg{classes}, and \mg{methods} for grouped
      data in \splus
    \item \spc \mg{Trellis} displays for grouped data
    \end{itemize}
  \end{enumerate}
\end{slide*}
\begin{slide*}
  \Ti{Outline (cont.'d)}
  \begin{enumerate}
  \item[3.] \rd{Linear mixed-effects models}
    \begin{itemize}
    \item brief review of linear models
    \item linear mixed-effects (\mg{LME}) model
    \item \spc \mg{Multilevel} LME models
    \item the \mg{lme} function in \splus
    \end{itemize}
  \item[4.] \rd{Extending the basic LME model}
    \begin{itemize}
    \item modelling the random effects \mg{covariance} structure
    \item variance functions for modelling \mg{heteroscedasticity}
    \item correlation structures: \mg{serial}, \mg{spatial}
    \item  S-PLUS capabilities for extending and customizing the
      \mg{lme} function
    \item the \mg{gls} function in \splus
    \end{itemize}
  \end{enumerate}
\end{slide*}
\begin{slide*}
  \Ti{Outline (cont.'d)}
  \begin{enumerate}
  \item[5.] \rd{Nonlinear mixed-effects models}
    \begin{itemize}
    \item \spc \mg{nonlinear} regression model
    \item starting estimates: \mg{self-starting} models
    \item nonlinear mixed-effects (\mg{NLME}) models
    \item the \mg{nlme} function in \splus
    \item extending the basic \mg{NLME} model
    \item using \mg{covariates} with \rd{nlme}
    \end{itemize}
  \end{enumerate}
\end{slide*}
\begin{slide*}
  \Ti{Some examples}

  {\red Orthodont} -- orthodontic growth data\\[2ex]
  Distance (in millimeters) between the pituitary and the
  pterygomaxillary fissure measured every two years for 27 children
  (16 males, 11 females) from age 8 until age 14. Reported in
\cite{pot64roy}
{\tiny
\begin{verbatim}
    distance age Subject  Sex 
  1     26.0   8     M01 Male
  2     25.0  10     M01 Male
  3     29.0  12     M01 Male
  4     31.0  14     M01 Male
  5     21.5   8     M02 Male
  6     22.5  10     M02 Male
. . .
104     19.5  14     F10 Female
105     24.5   8     F11 Female
106     25.0  10     F11 Female
107     28.0  12     F11 Female
108     28.0  14     F11 Female
\end{verbatim}}
  \begin{itemize}
  \item balanced data
  \item linear growth with time
\fig{Photo_4}{This is a test.  This is only a test.}
  \item similar patterns across individuals, but with different
    intercepts and slopes \cite{har98dev}
  \end{itemize}
\end{slide*}

\bs
\small
\bibliographystyle{abbrv}
\bibliography{feh}
\es

\bs{Abstract}
{\small
. . . . .
}
\es

\end{document}